\documentclass[10pt,a4paper]{article}

\usepackage[top=0.32in,bottom=0.46in,left=0.52in,right=0.52in]{geometry}
\usepackage[french]{babel}
\usepackage{fontspec}
\usepackage{xcolor}
\usepackage{hyperref}
\usepackage{graphicx}
\usepackage{enumitem}
\usepackage{tikz}
\usepackage{array}
\usepackage{microtype}
\usepackage{ragged2e}

\definecolor{headerbg}{HTML}{2F3440}
\definecolor{sectionfill}{HTML}{E9E9E9}
\definecolor{linegray}{HTML}{B7B7B7}
\definecolor{linkblue}{HTML}{0B57D0}

\hypersetup{
  colorlinks=true,
  urlcolor=linkblue,
  linkcolor=linkblue
}

\pagestyle{empty}
\setlength{\parindent}{0pt}
\setlength{\parskip}{0pt}
\linespread{1.02}
\pretolerance=10000
\tolerance=4000
\emergencystretch=3em
\hyphenpenalty=10000
\exhyphenpenalty=10000
\setlist[itemize]{
  label=\textbullet,
  leftmargin=1.15em,
  itemsep=0.10em,
  topsep=0.14em,
  parsep=0pt,
  partopsep=0pt
}

\defaultfontfeatures{Ligatures=TeX}
\setmainfont[
  Path = C:/Windows/Fonts/,
  UprightFont = segoeui.ttf,
  BoldFont = segoeuib.ttf,
  ItalicFont = segoeuii.ttf,
  BoldItalicFont = segoeuiz.ttf,
  Scale = 0.97
]{Segoe UI}

\newsavebox{\timelinebox}
\newlength{\timelineheight}

\newcommand{\sectionbadge}[1]{%
  \vspace{0.12em}%
  {\setlength{\fboxsep}{2.8pt}\colorbox{sectionfill}{\textbf{#1}}}\par\vspace{0.32em}
}

\newcommand{\headerlink}[2]{\href{#1}{\textcolor{white}{#2}}}

\newcommand{\contactcell}[1]{%
  \begin{minipage}[t]{0.23\linewidth}
    \centering\footnotesize\color{white}#1
  \end{minipage}%
}

\newcommand{\roundphoto}[1]{%
  \begin{tikzpicture}
    \clip (0,0) circle (0.92cm);
    \node at (0,0) {\includegraphics[width=1.84cm,height=1.84cm]{#1}};
    \draw[linegray,line width=0.6pt] (0,0) circle (0.92cm);
  \end{tikzpicture}%
}

\newcommand{\timelineentry}[5]{%
  \sbox{\timelinebox}{%
    \parbox[t]{0.93\linewidth}{%
      \RaggedRight
      \textbf{#1}\par
      {\small #2 \hfill \textit{#4}}\par
      {\small\textit{#3}}\par
      \vspace{0.05em}
      \begin{itemize}
        #5
      \end{itemize}
    }%
  }%
  \setlength{\timelineheight}{\dimexpr\ht\timelinebox+\dp\timelinebox\relax}%
  \noindent
  \begin{minipage}[t]{0.045\linewidth}
    \vspace{0pt}%
    \centering
    \begin{tikzpicture}[baseline=(dot.center)]
      \pgfmathsetlengthmacro{\timelineend}{\timelineheight-0.55em}
      \coordinate (top) at (0,0.72em);
      \coordinate (dot) at (0,0.14em);
      \coordinate (bottom) at (0,-\timelineend);
      \draw[linegray,line width=0.5pt] (top) -- (bottom);
      \draw[linegray,fill=white,line width=0.5pt] (dot) circle (0.078);
    \end{tikzpicture}%
  \end{minipage}%
  \hspace{0.01\linewidth}%
  \usebox{\timelinebox}\par\vspace{0.30em}
}

\newcommand{\sidebarentry}[3]{%
  \RaggedRight
  \textbf{#1}\par
  {\small #2}\par
  {\small\textit{#3}}\par
  \vspace{0.32em}
}

\newcommand{\skillentry}[2]{%
  \RaggedRight
  \textbf{#1}\par
  {\small #2}\par
  \vspace{0.32em}
}

\newcommand{\cvheader}{%
  \noindent
  \begin{minipage}[c][1.95cm][c]{0.12\textwidth}
    \centering
    \roundphoto{cv-oussama-photo.png}
  \end{minipage}
  \hfill
  \begin{minipage}[c][1.95cm][c]{0.84\textwidth}
    \raggedright
    {\fontsize{21.5}{23.5}\selectfont\textbf{OUSSAMA GOURIDECHE}}\par
  \end{minipage}

  \vspace{0.18em}

  {\setlength{\fboxsep}{8pt}%
  \colorbox{headerbg}{%
    \parbox{\dimexpr\textwidth-2\fboxsep\relax}{%
      \contactcell{\headerlink{mailto:oussgourideche@gmail.com}{oussgourideche@gmail.com}}
      \hfill
      \contactcell{Taza, Maroc}
      \hfill
      \contactcell{+212 6 60 50 80 66}
      \hfill
      \contactcell{Né le 05/05/1995 \quad Permis B}
    }%
  }}
}

\begin{document}

\cvheader

\vspace{0.16em}
{\fontsize{9.9}{11.55}\selectfont

\noindent
\begin{minipage}[t]{0.64\textwidth}
  \sectionbadge{Expériences professionnelles}

  \timelineentry
    {Ingénieur opérationnel}
    {IQRA Maroc}
    {2025}
    {Casablanca, Maroc}
    {
      \item Pilotage d'équipes terrain avec planification de charge, répartition des ressources et suivi opérationnel quotidien.
      \item Coordination interservices, gestion des aléas et analyse des KPIs pour soutenir l'amélioration continue.
    }

  \timelineentry
    {Chef de projet intégrateur de systèmes d'automatisation des opérations logistiques}
    {ALINA distribution}
    {09/2024}
    {Maroc}
    {
      \item Pilotage du projet de mise en place d'un entrepôt logistique avec structuration des flux et amélioration continue de l'exploitation.
      \item Mise en place des KPIs, suivi des performances et affectation des équipes selon les besoins opérationnels.
      \item Intégration d'un système ERP/WMS, suivi des stocks et optimisation de la traçabilité des opérations.
    }

  \timelineentry
    {Responsable préparation de commandes et contrôle marchandise}
    {ID Logistique (La Voie Express)}
    {09/2023 -- 08/2024}
    {Skhirat, Maroc}
    {
      \item Supervision de la préparation et du contrôle marchandise avec suivi des niveaux de service, de la qualité et de la conformité client.
      \item Suivi des consommables, contrôle d'exécution des clauses contractuelles et reporting KPI de l'activité plateforme.
    }

  \timelineentry
    {Manager opérationnel et commercial}
    {Myopla}
    {08/2022 -- 09/2023}
    {Maroc}
    {
      \item Encadrement d'une équipe B2C avec planification de production, suivi des objectifs et coordination de l'activité.
      \item Suivi qualité, satisfaction client, recrutement et montée en compétence des collaborateurs.
    }

  \timelineentry
    {Responsable contrôle et planification logistique}
    {SNCF}
    {08/2021 -- 08/2022}
    {France}
    {
      \item Planification des tournées taxi, suivi GPS des chauffeurs et coordination des trajets selon les contraintes opérationnelles.
      \item Facturation sous DOLIBARR, suivi de la consommation gasoil et appui au pilotage de l'équipe.
    }

  \timelineentry
    {Manager de la chaîne logistique (intégrateur du système d'automatisation et digitalisation de flux)}
    {AFD Tech}
    {2019 -- 2022}
    {Paris, France}
    {
      \item Gestion des opérations de stockage et d'approvisionnement avec coordination achats et négociation fournisseurs.
      \item Supervision des inventaires, contrôle des flux et suivi des besoins logistiques rattachés à l'exploitation.
      \item Intégration du système de gestion logistique avec participation à la mise en place d'un système développé de gestion logistique.
    }
\end{minipage}
\hfill
\begin{minipage}[t]{0.32\textwidth}
  \sectionbadge{Diplômes et Formations}

  \sidebarentry
    {Formation IT et sécurité réseau}
    {AFD University}
    {2021}

  \sidebarentry
    {Master manager de la chaîne logistique}
    {CNAM, RNCP 7 -- obtenu}
    {2019 -- 2021 \textbar{} Paris}

  \sidebarentry
    {Licence ST génie logistique}
    {UVSQ}
    {France}

  \sidebarentry
    {DUT logistique et transport}
    {EST}
    {2013 -- 2015 \textbar{} Oujda}

  \sidebarentry
    {Baccalauréat STG}
    {Diplôme obtenu}
    {}

  \sectionbadge{Compétences techniques}

  \skillentry
    {Outils métiers}
    {ERP, WMS, DOLIBARR, suivi GPS, tableaux de bord KPI}

  \skillentry
    {Flux logistiques}
    {réception, stockage, inventaires, préparation de commandes, distribution}

  \skillentry
    {Pilotage d'activité}
    {planification, coordination terrain, allocation des ressources, suivi qualité}

  \skillentry
    {Développement et data}
    {Java, SQL, HTML, CSS, JavaScript, modélisation de données et requêtes}
\end{minipage}

}

\newpage

{\fontsize{9.9}{11.55}\selectfont

\noindent
\begin{minipage}[t]{0.64\textwidth}
  \sectionbadge{Expériences professionnelles (suite)}

  \timelineentry
    {Gestionnaire de stock}
    {UPS France}
    {2018 -- 2019}
    {France}
    {
      \item Gestion des entrées et sorties de stock, contrôle des références et suivi d'inventaire.
      \item Fiabilisation des niveaux de stock, préparation des inventaires et suivi opérationnel des mouvements logistiques.
    }
\end{minipage}
\hfill
\begin{minipage}[t]{0.32\textwidth}
  \sectionbadge{Compétences techniques (suite)}

  \skillentry
    {Analyse et amélioration}
    {suivi de performance, lecture des écarts, mise en place de plans d'action, optimisation des opérations}

  \skillentry
    {Culture IT}
    {fondamentaux réseaux, sécurité opérationnelle, outils numériques métier, ERP/WMS et environnement bureautique}
\end{minipage}

\vspace{0.34em}

\sectionbadge{Compétences complémentaires}
\begin{itemize}
  \item Suivi et analyse des indicateurs de performance pour piloter l'activité, identifier les écarts et prioriser les actions correctives.
  \item Structuration des flux de réception, stockage, préparation et distribution avec attention portée à la traçabilité et à la qualité d'exécution.
  \item Coordination d'équipes terrain et communication interservices pour fluidifier l'exploitation et sécuriser la continuité opérationnelle.
  \item Utilisation d'outils numériques métier pour le reporting, la facturation, le suivi logistique et la consolidation d'informations d'exploitation.
\end{itemize}

}

\end{document}
